\chapter{优化方法与算法设计}

\section{系统总体设计}

\subsection{设计目标}

围绕移动端异构多核 CPU 上的大模型推理过程，调度优化面临两个直接约束。其一，Prefill 与 Decode 的负载特征差异明显，难以共用同一套线程数和绑核配置。其二，移动端运行环境存在线程迁移、背景任务抢占和核心能力不均等问题，单纯依赖固定参数的静态划分难以长期保持稳定。为此，本章将优化目标归纳为三点：一是针对不同阶段分别选择更合适的线程数与核心集合；二是利用真实性能比构造更合理的初始任务分配；三是在 Prefill 阶段引入任务窃取机制，减轻尾部阻塞。

方法设计遵循最小侵入原则。换言之，不改变原有推理流程的基本组织方式，只在任务划分和并行执行入口处引入阶段感知的调度控制。这样做的好处在于，优化过程能够直接落到现有运行环境中，后续实验也更容易复现。

\subsection{总体思路}

整体流程如图~\ref{fig:chapter3_overview_placeholder} 所示。推理开始前，先根据目标机型完成线程数和绑核参数的搜索，并通过一组可复用的配置描述 Prefill 与 Decode 的执行状态。进入运行阶段后，Prefill 先根据核心相对能力生成初始任务区间，再通过无锁双端队列 work stealing 机制完成负载均衡；Decode 则采用更轻量的动态领取方式，以避免短步迭代中额外调度开销过大。

\begin{figure}[ht]
  \centering
  \fbox{
    \parbox{0.9\textwidth}{
      \vspace{1.2cm}
      图位预留：系统总体流程图。\newline
      建议内容：从左到右画出 ``阶段参数搜索''、``Prefill 执行路径''、``Decode 执行路径'' 三个模块。\newline
      Prefill 模块中画出 ``真实性能比 -> 初始任务区间 -> 本地执行 -> work stealing''。\newline
      Decode 模块中画出 ``均分区间 -> 原子领取 -> 执行完成''。\newline
      图中需要标明阶段切换点、绑核参数输入和日志统计输出。
      \vspace{1.2cm}
    }
  }
  \caption{阶段感知调度总体流程示意图}
  \label{fig:chapter3_overview_placeholder}
\end{figure}

\subsection{设计原则}

这一调度方案遵循三条原则。

第一，阶段分离。Prefill 关注总工作量较大的长任务处理，关键问题是负载均衡与尾部线程拖慢整体结束时间；Decode 关注短步迭代和连续响应，关键问题是调度开销和时延抖动。因此，两阶段必须分别配置线程数和核心集合。

第二，先给出合理初值，再做运行期修正。异构多核上的任务划分不能完全依赖运行期竞争，否则调度成本会上升；也不能完全依赖静态边界，否则一旦出现扰动就容易失衡。更合适的办法是先利用真实性能比构造较合理的初始区间，再在运行中通过 stealing 修正剩余不均衡。

第三，重方法、轻机制。第 3 章重点说明方法设计与算法逻辑，而不是罗列大量实现接口。与章节论证直接相关的，主要是阶段划分规则、任务组织方式、窃取规则和调度统计项。

\section{分阶段线程数与绑核配置}

\subsection{问题分析}

在同一台手机上，Prefill 与 Decode 的最优线程数通常并不一致。Prefill 阶段计算量更大，对大核利用更敏感，线程数偏少会导致高性能核心未被充分使用；线程数过多又会引入同步、访存争用和额外调度开销。Decode 阶段则不同，它以短步迭代为主，单步工作量较小，若并发度设置过高，调度开销往往先于计算收益出现。

这意味着，统一线程数与统一 affinity 的做法并不稳妥。若所有阶段共用同一组参数，通常只能得到折中解，而不是各阶段的较优解。因而，线程数与绑核配置需要按阶段独立确定。

\subsection{配置方法}

配置过程分为两步。第一步确定候选核心集合。对于 Prefill，优先保留高性能核心；对于 Decode，则在保证性能稳定的前提下避免过多线程竞争。第二步在给定核心集合上搜索线程数。搜索的判断依据不是单次最好值，而是重复实验下的平均表现和波动情况。若某一线程数配置虽然偶尔较快，但波动较大，则不宜作为最终配置。

从方法上看，这一配置过程并不依赖复杂优化器，而是采用可解释的启发式搜索。这样做的原因是，移动端机型差异较大，过度复杂的离线模型未必能带来更稳定的实际收益。相反，围绕目标机型直接完成参数搜索，更贴近毕业论文需要强调的工程可复现路径。

\subsection{算法描述}

算法~\ref{alg:phase_config_search} 给出了分阶段线程数与绑核搜索的基本流程。其输入为候选核心集合、线程数上界和测试负载；输出为 Prefill 与 Decode 各自采用的线程数与核心集合。

\begin{algorithm}[ht]
\caption{分阶段线程数与绑核参数搜索}
\label{alg:phase_config_search}
\KwIn{候选核心集合 $C$，线程数搜索范围 $T$，测试负载 $W$}
\KwOut{Prefill 参数 $(C_p, t_p)$，Decode 参数 $(C_d, t_d)$}
根据核心能力对 $C$ 排序\;
为 Prefill 选择高性能核心候选集合\;
为 Decode 选择稳定性更好的核心候选集合\;
\ForEach{候选线程数 $t \in T$}{
  在 Prefill 负载上测试 $(C_p, t)$ 的平均时延与波动\;
  记录较优的 $(C_p, t_p)$\;
}
\ForEach{候选线程数 $t \in T$}{
  在 Decode 负载上测试 $(C_d, t)$ 的平均时延与波动\;
  记录较优的 $(C_d, t_d)$\;
}
返回 $(C_p, t_p)$ 和 $(C_d, t_d)$\;
\end{algorithm}

\subsection{图示建议}

这一部分建议配两张图。

第一张图放在本节后，画 Prefill 与 Decode 随线程数变化的性能曲线。横轴为线程数，纵轴可用时间或吞吐。图中应清楚看到两条曲线的最优点不重合。

第二张图可以画绑核集合示意图。建议用手机 CPU 拓扑图表示 big、mid、little 三类核心，并在图中标出 Prefill 与 Decode 分别使用的核心集合。

\section{Prefill 阶段的初始任务组织}

\subsection{设计动机}

Prefill 的总任务量较大，若一开始就简单均分，异构核心之间的能力差异会很快转化为执行时间差异。高性能核心可能先完成任务并空闲，而低性能核心仍保留大量工作。此时即使引入 stealing，窃取线程也需要在后期承担更多补偿性工作，整体效率并不理想。

因此，Prefill 阶段的第一步不是直接窃取，而是先构造更合理的本地任务区间。这里所说的“更合理”，指的是初始工作量与核心处理能力大致匹配。能力更强的核心应拿到更大的初始区间，能力较弱的核心应拿到更小的区间。这样，stealing 的角色就从“替代主划分”变成“修正残余不均衡”。

\subsection{真实性能比的作用}

构造初始任务区间需要一个能够刻画核心相对能力的指标。固定经验参数虽然实现简单，但在机型变化、模型变化和系统负载变化时容易失准。更稳妥的做法是通过实测获得各核心或各簇的相对性能比，再按照这一比值分配初始工作量\cite{wang2021asymo,abdelgawad2022bertperf}。

这里不追求精细的体系结构建模，而是把性能比作为任务组织的工作参数。只要该参数能够较稳定地反映不同核心承担工作量的合理比例，就足以支撑初始区间划分。

\subsection{初始区间生成规则}

设总任务量为 $N$，活动线程数为 $P$，第 $i$ 个线程对应核心的相对能力为 $w_i$，满足
\[
\sum_{i=1}^{P} w_i = 1.
\]
则第 $i$ 个线程的目标初始任务量可表示为
\[
n_i = N \cdot w_i.
\]
在实际划分时，由于任务边界必须取整数，故采用累积划分方式生成各线程边界。这样可以避免局部舍入误差不断累积，保证所有初始区间首尾相接、总和不超过任务总量。

\subsection{算法描述}

算法~\ref{alg:weighted_seed} 给出了基于性能比的初始区间生成方法。其输出不是最终调度结果，而是每个线程在 work stealing 开始前的本地任务种子区间。

\begin{algorithm}[ht]
\caption{基于真实性能比的初始区间生成}
\label{alg:weighted_seed}
\KwIn{总任务量 $N$，活动线程数 $P$，性能比 $\{w_1,w_2,\ldots,w_P\}$}
\KwOut{各线程初始区间边界 $\{b_0,b_1,\ldots,b_P\}$}
令 $b_0 \leftarrow 0$\;
令累计量 $s \leftarrow 0$\;
\For{$i \leftarrow 1$ \KwTo $P$}{
  $s \leftarrow s + N \cdot w_i$\;
  计算边界 $b_i \leftarrow \mathrm{round}(s)$\;
  若 $b_i < b_{i-1}$，则令 $b_i \leftarrow b_{i-1}$\;
  若 $b_i > N$，则令 $b_i \leftarrow N$\;
}
返回边界序列 $\{b_0,b_1,\ldots,b_P\}$\;
\end{algorithm}

\begin{figure}[ht]
  \centering
  \fbox{
    \parbox{0.9\textwidth}{
      \vspace{1.0cm}
      图位预留：Prefill 初始区间生成示意图。\newline
      建议内容：横向条带表示总任务区间，按不同颜色切分为若干线程的初始本地区间。\newline
      上半部分画 ``均分''，下半部分画 ``按真实性能比加权划分''。\newline
      图中标出大核对应区间更长，小核对应区间更短。
      \vspace{1.0cm}
    }
  }
  \caption{Prefill 初始任务区间组织示意图}
  \label{fig:prefill_seed_placeholder}
\end{figure}

\section{Prefill 阶段的 work stealing 机制}

\subsection{基本思想}

当初始区间建立后，每个线程先处理自己的本地任务。当本地任务耗尽时，再尝试从其他线程剩余任务中窃取一部分继续执行。这样做的目的是让空闲线程重新参与计算，减轻最慢线程对阶段结束时间的拖累\cite{blumofe1999workstealing,chase2005deque}。

Prefill 采用双端队列式的任务组织方式。拥有者线程优先从本地一端连续取任务；窃取线程则从受害者的另一端取走剩余区间的一部分。这样的安排有两个好处。第一，本地线程和窃取线程在访问方向上尽量分离，冲突较小。第二，本地线程能够连续处理相邻任务块，局部性较好。

\subsection{本地优先原则}

work stealing 并不意味着所有线程从一开始就频繁竞争共享任务。更合适的执行方式是本地优先，即线程先尽量消化本地任务，再在必要时发起窃取。这样可以避免过早进入跨线程竞争，也能更充分利用初始区间中已经体现出来的异构能力差异。

在 Prefill 阶段，本地优先原则尤其重要。若本地任务尚未执行完就频繁窃取，不仅会增加同步冲突，还会破坏原本按核心能力构造的初始分配。

\subsection{窃取规则}

窃取时并不一次性取走全部剩余任务，而是取走受害者剩余区间的一部分。这样做是为了避免“窃取得过多，新的长尾又出现”的情况。较稳妥的办法是从受害者剩余任务的尾部取走一半或接近一半的工作量，使拥有者仍保留继续推进的能力，窃取者则得到足够的补偿任务。

受害者选择采用固定顺序轮询方式即可，不必引入复杂的随机化机制。对移动端推理而言，调度路径越短越好。只要能够在较低成本下找到仍有剩余任务的线程，方法就已经足够有效。

\subsection{局部块大小调整}

异构多核上的 work stealing 还有一个容易忽略的问题：即使初始区间已经按性能比划分，各线程在本地推进时也不应使用完全相同的块大小。高性能核心可以采用稍大的本地块，以减少频繁领取任务造成的额外开销；低性能核心则可采用相对保守的块大小，避免块粒度过大引发新的拖尾。

因此，本地块大小并不是全局固定常数，而是依据初始区间大小做比例调整。直观地说，初始区间越大，说明该线程承担的基础工作越多，本地推进步长也可相应放大；初始区间越小，本地推进步长则保持较小值。

\subsection{算法描述}

算法~\ref{alg:prefill_ws} 描述了 Prefill 阶段的基本执行逻辑。为便于阅读，伪代码省略了底层原子操作细节，仅保留方法层面的关键步骤。

\begin{algorithm}[ht]
\caption{Prefill 阶段的本地优先 work stealing}
\label{alg:prefill_ws}
\KwIn{各线程初始区间 $\{Q_1,Q_2,\ldots,Q_P\}$}
\KwOut{Prefill 阶段全部任务执行完成}
\ForPar{每个线程 $i$}{
  \While{$Q_i$ 非空}{
    从 $Q_i$ 本地领取一个任务块并执行\;
  }
  \While{仍存在非空队列}{
    选择一个受害者线程 $j$\;
    \If{$Q_j$ 非空}{
      从 $Q_j$ 的尾部窃取一部分任务\;
      执行窃取得到的任务块\;
    }
    \Else{
      尝试下一个受害者线程\;
    }
  }
}
\end{algorithm}

\begin{figure}[ht]
  \centering
  \fbox{
    \parbox{0.9\textwidth}{
      \vspace{1.0cm}
      图位预留：Prefill 阶段 work stealing 示意图。\newline
      建议内容：画四个线程的任务队列。前两步显示 ``本地执行''，后两步显示某个线程本地任务耗尽后，从另一线程队列尾部取走一部分任务。\newline
      图中最好标注 ``owner 从前端推进''、``thief 从尾部窃取''，以及窃取前后队列长度变化。
      \vspace{1.0cm}
    }
  }
  \caption{Prefill 阶段任务窃取过程示意图}
  \label{fig:prefill_ws_placeholder}
\end{figure}

\subsection{统计量设计}

仅给出总时延还不足以解释 work stealing 的作用，因此还需要记录若干调度统计量。较重要的指标包括：本地任务领取次数、本地领取成功次数、窃取尝试次数、窃取成功次数、窃取失败次数、被窃取任务量以及各线程最终承担的任务总量。通过这些指标，可以判断性能提升究竟来自更合理的初始分配，还是来自运行期长尾缓解。

\section{Decode 阶段的轻量动态调度}

\subsection{设计原因}

Decode 阶段与 Prefill 不同，任务短、迭代频繁。若在这一阶段仍采用多队列 stealing，线程之间会反复进行受害者选择和任务竞争，调度成本会迅速放大。因此，Decode 更适合使用轻量动态领取方式，而不是复用 Prefill 的 work stealing 机制。

\subsection{调度方式}

Decode 的基本做法是将总任务区间均匀展开，并设置一个共享游标。线程每次通过原子方式领取一个连续小区间，执行完成后再继续领取，直到总任务耗尽。这种方式只有一个共享入口，路径短、实现简单，适合短步迭代场景。

\subsection{算法描述}

算法~\ref{alg:decode_dynamic} 给出了 Decode 阶段的动态领取过程。

\begin{algorithm}[ht]
\caption{Decode 阶段的动态任务领取}
\label{alg:decode_dynamic}
\KwIn{总任务量 $N$，块大小 $s$}
\KwOut{Decode 阶段全部任务执行完成}
令共享游标 $c \leftarrow 0$\;
\ForPar{每个线程}{
  \While{true}{
    原子领取区间起点 $l \leftarrow c$，并令 $c \leftarrow c + s$\;
    \If{$l \geq N$}{
      退出循环\;
    }
    令 $r \leftarrow \min(N, l + s)$\;
    执行区间 $[l, r)$ 内的任务\;
  }
}
\end{algorithm}

\subsection{与 Prefill 调度的区别}

两阶段的差异可以概括为一句话：Prefill 关注长尾均衡，Decode 关注低开销稳定。前者需要在本地优先的基础上保留窃取能力，后者则应尽量缩短任务领取路径。也正因如此，统一调度策略很难同时兼顾两个阶段。

\section{本章小结}

本章围绕阶段感知调度的核心方法展开了说明。首先，线程数与绑核参数按阶段独立配置，避免统一参数带来的折中损失。其次，Prefill 阶段利用真实性能比生成初始任务区间，并在此基础上采用本地优先的 work stealing 机制缓解尾部阻塞。再次，Decode 阶段采用轻量动态领取方式，以控制短步迭代中的调度成本。第 4 章将在真实手机平台上对这些方法进行实验评估，并补充相应的曲线图、示意图和性能结果。
